\documentclass[10pt,a4paper]{amsart}
\usepackage[margin=19mm]{geometry}
\usepackage{amsmath,amssymb,mathtools}
\usepackage[scr=rsfs,cal=euler]{mathalfa}
\usepackage{xfrac}
\usepackage[unicode,hidelinks,bookmarks=false]{hyperref}
\newcommand{\ie}{i.e.\ }
\newcommand{\cf}{cf.\ }
\newcommand{\resp}{respectively\ }
\newcommand{\Subsection}{\subsection}
\providecommand{\nobreakdash}{\nobreak\hbox{-}}
\setlength{\emergencystretch}{2em}
\multlinegap0pt
\usepackage{amssymb}
\usepackage{tikz}
\usepackage{xcolor}
\newtheorem{proposition}{Proposition}
\newcommand{\OrdP}{\mathrm{Ord}_{P}}
\title{The Parry Order of Perron Numbers}
\author{Kevin G Hare, Hachem Hichri}
\date{Source: arXiv:2603.19554v1 (CC BY 4.0)}
\begin{document}
\maketitle

\noindent\emph{Source and licence.} The Parry Order of Perron Numbers. Version arXiv:2603.19554v1 (CC BY 4.0), distributed by arXiv under the Creative Commons Attribution 4.0 International licence, http://creativecommons.org/licenses/by/4.0/, as recorded in the arXiv licence metadata for this exact version and shown on the abstract page https://arxiv.org/abs/2603.19554v1. Full licence notice: \texttt{licenses/arxiv-2603.19554-CC-BY-4.0.txt}.\par\smallskip

\noindent\textit{Editorial scope.} Counted unit: the article's proposition proposition (source0.tex lines 661--663). The article's complete original proof of it is delivered with it (source0.tex lines 666--697, one \texttt{proof} environment of 32 lines). No mathematics is written, repaired or re-derived: the statement and the proof are the article's own lines, in the article's own notation, and no retained line is substituted or re-flowed.  No further source-local context is needed: every cross-reference of every retained line resolves inside the two delivered spans.  Nothing was omitted from any retained span, so the package carries no omission record.  No retained line contains a citation, so the wrapper carries no bibliography and no bibliography entry.  No retained line contains a figure environment, an image-inclusion command or drawing code, so no image is delivered and the package declares no asset dependency.  Editorial formatting only, in addition: an AMS \texttt{amsart} preamble that restates the article's own declarations and macros used by the retained lines, namely the environments \texttt{proposition} on separate counters, together with the standard packages those lines need.  The retained environments print in this document as Proposition 1 in order of appearance, whereas the article prints the same items with the numbering of its own full document.  The complete native source of the article as distributed in the arXiv e-print for this exact version is bundled unchanged as \texttt{sources/196784/source0.tex}.
\begin{proposition}
    For all $N$ there exists a non-Pisot Parry number $\beta$ such that $\OrdP(\beta) > N$.
\end{proposition}

\begin{proof}
    Let $\gamma(d), \alpha(d)$ and $\beta(d)$ be the root of $x^3-d x^2-2 x + d$ where
    $\gamma(d) < -1 < \alpha(d) < 1 < \beta(d)$.

    We can show that
    \begin{align}
    \gamma(d) & = -1-\frac{1}{2d} + \frac{3}{8 d^2} - \frac{55}{128 d^4} + \dots, \label{eq:1} \\
    \alpha(d) & = 1-\frac{1}{2d} - \frac{3}{8 d^2} + \frac{55}{128 d^4} + \dots, \label{eq:2} \\
    \beta(d) & = d + \frac{1}{d} - \frac{1}{d^5} + \frac{2}{d^7} - \frac{1}{d^9} + \dots.  \label{eq:3}
    \end{align}

    We observe from this that for sufficiently large $d$ that $\beta_d$ is not a Pisot number.  Hence it's Parry order is bounded for any fixed $d$.

    Letting $p_n(x) = (x-\beta(d)^n) (x-\alpha(d)^n) ( x- \gamma(d)^n) = x^3 - a_n(d) x^2 - b_n(d) x + c_n(d)$.

    Using equations \eqref{eq:1}, \eqref{eq:2} and \eqref{eq:3} we get 
    \begin{align*}
        a_n(d) & = 
            \gamma(d)^n + \beta(d)^n + \alpha(d)^n, \\
        b_n(d) & = -(\gamma(d)^n \beta(d)^n + \gamma(d)^n \alpha(d)^n + \alpha(d)^n \beta(d)^n). \\
    \end{align*}

    Let $d$ be odd, and consider $a_d(d^2)$ and $b_d(d^2)$. 
    This simplify to 
    \begin{align*}
        a_d(d^2) & = d^{2d} \left( 1 + \frac{1}{d^2} + \frac{1}{2 d^6} - \frac{1}{2 d^7} + \frac{1}{6 d^9} + \dots \right), \\
        b_d(d^2) & = d^{2d} \left(-\frac{1}{d} - \frac{1}{24 d^3} - \frac{1}{2 d^4} - \frac{881}{1920 d^5} - \frac{1}{48 d^6} - \dots \right).         
    \end{align*}
    In particular this implies that for large odd $d$ that we have $\OrdP(\beta_{d^2}) \geq d$.

This proves the desired result.
\end{proof}

\end{document}
